% Created 2025-06-25 Wed 18:49
% Intended LaTeX compiler: pdflatex
\documentclass[11pt]{article}
\usepackage[utf8]{inputenc}
\usepackage[T1]{fontenc}
\usepackage{graphicx}
\usepackage{longtable}
\usepackage{wrapfig}
\usepackage{rotating}
\usepackage[normalem]{ulem}
\usepackage{amsmath}
\usepackage{amssymb}
\usepackage{capt-of}
\usepackage{hyperref}
\author{fcalado}
\date{\today}
\title{}
\hypersetup{
 pdfauthor={fcalado},
 pdftitle={},
 pdfkeywords={},
 pdfsubject={},
 pdfcreator={Emacs 29.3 (Org mode 9.6.15)}, 
 pdflang={English}}
\begin{document}

\tableofcontents

Final Project: How To Build A Bot Tutorial

Assignment description:
\begin{itemize}
\item choose a social media app, like instagram, tiktok, linkedin, or
another app of your choice.
\item research some tutorials for scraping and/or creating a bot for that
app. Make sure the tutorial is recent (in the last year, at
minimum).
\item with a partner or by yourself, create a tutorial that you will use
to teach your classmates how to scrape or create a bot on that app.
\end{itemize}

Deliverables:
\begin{itemize}
\item written tutorial, in markdown (\texttt{.md}) format:
\begin{itemize}
\item with each step described clearly, and code blocks to include code
examples.
\item see the `day8\_twitter\_bot.md` file in /class\_notes folder for a
sample of an \texttt{.md} file.
\end{itemize}
\item vitual presentation, 20 - 30 minutes:
\begin{itemize}
\item you will present the tutorial like a lesson, where you walk your
classmates through the process of using the tool.
\end{itemize}
\end{itemize}

Examples:
\begin{itemize}
\item Yin, Piotr Sapiezynski and Leon. \href{https://inspectelement.org/browser\_automation.html}{Browser Automation}, Blog.
\item Instagrapi, \href{https://www.youtube.com/watch?v=cW7kMeOUr20}{instagrapi tutorial}, Youtube.
\item \href{https://www.geeksforgeeks.org/make-an-instagram-bot-with-python/}{Make an Instagram Bot With Python}, Geeks for Geeks
\end{itemize}
\end{document}
